\section{What the triplet premise still guarantees in any dimension}\label{sec:orderedconductances}
The counterexample invalidates the proposed exclusion of ordinary cooling.
A substantial part of the exclusion principle nevertheless survives, under
only the $G$ inequalities and for any number of states.

Let $q_i=\pi_i(\beta_b)>0$ and define the symmetric reduced conductances
\[
k_{ij}=\frac{M_{ij}}{q_i}=\frac{M_{ji}}{q_j},\qquad i\ne j.
\]
These are conductances divided by both equilibrium weights. For an ordered
triplet $i<j<k$, its original-rate expression is
$G_{ijk}=a_{ik}q_j/q_i-a_{jk}=q_j(k_{ik}-k_{jk})$. Thus the condition
\begin{equation}\label{eq:orderedk}
k_{ik}\leq k_{jk}\quad\text{for all }i<j<k
\end{equation}
is exactly $G_{ijk}\leq0$ for every triplet. Failure of every original
predicate implies this condition, by the branch-equivalence lemma. No
$F$ inequality is used in the following theorem.

\begin{theorem}[No strong effect and no anomalous heating]\label{thm:orderedk}
Consider an irreducible reversible $n$-state generator, $n\geq2$, with distinct ordered
energies and $0<\beta_b<\infty$. Suppose \eqref{eq:orderedk} holds.
Then:
\begin{enumerate}
\item No nonbath Gibbs preparation at a finite positive temperature has zero
projection onto the entire slowest nonstationary eigenspace. Hence there is
no strong cooling or strong heating effect, including when that eigenspace
is degenerate.
\item For any $\beta_f>\beta_n>\beta_b$ and every finite $t\geq0$,
\[
\sum_i\frac{[p_i(t;\beta_f)-q_i]^2}{q_i}
>
\sum_i\frac{[p_i(t;\beta_n)-q_i]^2}{q_i}.
\]
Thus heating has no finite-time weighted Euclidean reversal. It also has
no asymptotic relative-entropy reversal, for any slow multiplicity.
\item If the slow mode is simple, its thermal overlap can be oriented to be
strictly negative throughout cooling, strictly positive throughout heating,
and strictly increasing throughout heating. Consequently there is no
asymptotic heating effect in any fixed norm.
\end{enumerate}
These statements do not exclude ordinary cooling, and do not assert
finite-time relative-entropy or total-variation ordering.
\end{theorem}

\begin{proof}
Write $L=M^\top$, $Q_k=\sum_{i=1}^kq_i$, and define the weighted prefix
contrasts, for $k=1,\ldots,n-1$, by
\[
h^{(k)}_i=\begin{cases}
1,&i\leq k,\\
-Q_k/q_{k+1},&i=k+1,\\
0,&i>k+1.
\end{cases}
\]
They have zero $q$-mean and are mutually orthogonal in
$\langle f,g\rangle_q=\sum_iq_if_ig_i$. Their squared norms are
$\eta_k=Q_kQ_{k+1}/q_{k+1}$. Put $e^{(k)}=h^{(k)}/\sqrt{\eta_k}$.
This is a weighted Helmert basis of the entire nonstationary function space.

Detailed balance makes the matrix $A_{kl}=\langle e^{(k)},Le^{(l)}\rangle_q$
symmetric. For $k<l$, direct multiplication gives
\begin{align}\label{eq:helmertoffdiag}
\sqrt{\eta_k\eta_l}A_{kl}
={}&Q_{l+1}\sum_{i=1}^kq_i(k_{k+1,l+1}-k_{i,l+1})\\
&+\sum_{j=l+2}^nq_j\sum_{i=1}^kq_i(k_{k+1,j}-k_{ij})\geq0.\nonumber
\end{align}
To check the identity, note that for every $i\leq k+1\leq l$,
$(Lh^{(l)})_i=-Q_{l+1}k_{i,l+1}-\sum_{j=l+2}^nq_jk_{ij}$.
The weighted contrast subtracts its value at $i=k+1$ from the weighted
prefix values. Each resulting difference is nonnegative by
\eqref{eq:orderedk}. Thus $A$ is a symmetric Metzler matrix. In particular
$e^{tA}$ is entrywise nonnegative for $t\geq0$, with positive diagonal.
This follows by choosing $c$ so that $A+cI$ is entrywise nonnegative and
using the exponential power series.

Let $g_k(\beta)=\sum_i e^{(k)}_i\pi_i(\beta)$. Apart from its positive
normalization $\sqrt{\eta_k}$, its numerator is
\begin{align*}
N_k(\beta)
&=\sum_{i=1}^k e^{-\beta E_i}-\frac{Q_k}{q_{k+1}}e^{-\beta E_{k+1}}\\
&=e^{-\beta E_{k+1}}\sum_{i=1}^k
\left[e^{\beta(E_{k+1}-E_i)}-e^{\beta_b(E_{k+1}-E_i)}\right].
\end{align*}
Every gap in this sum is strictly positive. Therefore every coordinate
$g_k$ is strictly negative for $\beta<\beta_b$, zero at the bath, and
strictly positive for $\beta>\beta_b$.

Each $g_k$ is also strictly increasing on the heating side. Indeed, put
\[
T_k=(1+Q_k/q_{k+1})e^{-\beta E_{k+1}}+
\sum_{j=k+2}^ne^{-\beta E_j},\qquad Z=N_k+T_k.
\]
For $\beta>\beta_b$, both $N_k,T_k$ are positive and
\[
\frac{N_k'}{N_k}>-E_{k+1},\qquad
\frac{T_k'}{T_k}\leq-E_{k+1}.
\]
The first inequality follows by differentiating the positive bracketed sum
in $N_k$; the second is its interpretation as the negative weighted average
of energies at least $E_{k+1}$. Consequently
$(N_k/Z)'=(N_k'T_k-N_kT_k')/Z^2>0$.

The Gibbs likelihood ratio expands as
$\pi(\beta)/q=\mathbf1+\sum_k g_k(\beta)e^{(k)}$.
Under the reversible evolution its nonconstant coordinate vector becomes
$e^{tA}g(\beta)$. For two heating preparations farther and nearer the bath,
$g(\beta_f)>g(\beta_n)>0$ entrywise. The nonnegative matrix $e^{tA}$ with
positive diagonal preserves both strict inequalities. Taking squared
Euclidean norms of the coordinate vectors proves the finite-time assertion.

For the slow-space conclusions let $\lambda_2$ be the largest eigenvalue
of $A$ and let $\Pi_2$ be its orthogonal spectral projector. Symmetry gives
\[
\Pi_2=\lim_{t\to\infty}e^{-\lambda_2t}e^{tA}.
\]
Hence $\Pi_2$ is entrywise nonnegative and nonzero. Since $g(\beta)$ has
strictly one sign in every coordinate, $\Pi_2g(\beta)\ne0$ for every
nonbath $\beta$. This proves absence of strong effects without assuming
simplicity. Moreover $\|\Pi_2g(\beta)\|^2$ is strictly increasing on the
heating side: its derivative is
$2(\Pi_2g)^\top(\Pi_2g')>0$. Strictness follows from any positive diagonal
entry of the nonzero nonnegative projector, since $g,g'>0$. The same
quantity controls the leading relative-entropy coefficient.

If the slow mode is simple, Perron theory for a symmetric Metzler matrix
allows its coordinate eigenvector $v$ to be chosen nonnegative and nonzero.
Then $b_2(\beta)=v^\top g(\beta)$ has the stated signs and heating
monotonicity. Its scale does not affect any ordering conclusion.
\end{proof}

\paragraph{A larger rate family.}
The proof in fact requires only that the right-hand side of
\eqref{eq:helmertoffdiag} be nonnegative for each $k<l$; the individual
conductance inequalities are sufficient, but not necessary. In four states
this larger, explicitly checkable family is given by just three inequalities:
\begin{align*}
Q_3(k_{23}-k_{13})+q_4(k_{24}-k_{14})&\geq0,\\
k_{24}-k_{14}&\geq0,\\
(q_1+q_2)k_{34}-q_1k_{14}-q_2k_{24}&\geq0.
\end{align*}
Every conclusion of Theorem~\ref{thm:orderedk} holds under these three
aggregate inequalities. More generally there are $(n-1)(n-2)/2$ linear
aggregate inequalities. They can be evaluated before any eigendecomposition.

For a simple four-state slow mode in this family, the anomalous cooling
region $\{\beta\in(0,\beta_b):\partial_\beta|b_2|>0\}$ is either empty or
one connected interval. Indeed $b_2<0$ on cooling and $b_2'>0$ throughout
heating, while the global derivative bound allows at most two stationary
points counting multiplicity. The negative set of $b_2'$ before the bath can therefore have at most
one connected component. In this larger family its left endpoint is not
determined by the present argument. Under the original full triplet premise,
Corollary~\ref{cor:tripletendpoint} proves that any such interval starts at
$\beta=0$ and supplies an exact endpoint criterion.

\paragraph{A corrected implication and its mechanism.}
For every finite $n$, failure of all the original triplet predicates excludes
strong effects and excludes asymptotic heating in relative entropy; with a
simple slow mode, heating is excluded in every fixed norm. The finite-time
weighted Euclidean statement is stronger for its specified distance. The
four-state counterexample therefore occupies a precise remaining possibility:
ordinary cooling without eliminating the slow mode. Its small thermal maximum
is compatible with a slow overlap that never changes sign on either side of
the bath.

The proof uses a different order structure from monotone energy eigenvectors.
A positive combination of prefix contrasts can have nonmonotone components,
so the counterexample does not conflict with this cone. Nevertheless its
weighted prefix sums are nonnegative. More directly, every Gibbs preparation
on one side of the bath has coordinates of the same sign in this basis; a
nonnegative slow spectral projector cannot cancel all of them. This is the
mechanism excluding strong effects. On heating, all those coordinates also
grow with the distance from the bath, which explains the additional exclusion.
The basis, Metzler positivity, and Perron theory are established tools; the
specific rate-to-contrast identity and its Mpemba consequences require a
separate priority assessment.
